% !PRINTED 647
\weberpage{647}
\begin{center}
\S.\,187.

\gesperrt{Reality Relations.}
\end{center}

If the field $\mathfrak{K}$ is a real one, then, as we have seen in §.\,180, there are two kinds of metacyclic equations of prime degree $n$: those with one real and $n-1$ imaginary roots, and those with all real roots. Likewise we have seen in §.\,165 that there are two kinds of cyclic equations of an even degree, namely those with all real and those with all imaginary roots.

% [sic] Print runs both the $k$ and $\tau$ lists through index $n-1$ here;
% the cyclic equation and radical system defined in §.185 run through $n-2$.
If now the cyclic equation whose roots are the quantities $k_0, k_1 \ldots k_{n-1}$ belongs to the first kind, that is, if $k_0, k_1 \ldots k_{n-1}$ are real and the radicals $\tau_0, \tau_1, \ldots \tau_{n-1}$ are also taken to be real, then formula (10) of §.\,186 shows that $\xi_0$ is real, and that $\xi_h$ and $\xi_{-h}$ are conjugate imaginary.

If on the other hand the $k$ are imaginary, then by §.\,165 $k_\nu$ and $k_{\nu+\frac{n-1}{2}}$ are conjugate imaginary, and by formula (8) in §.\,185 we can set, if $\Phi$ denotes a rational function,

\[
\sqrt[n]{R_{\nu+\frac{n-1}{2}}} = \Phi(k_\nu) \left(\sqrt[n]{R_\nu}\right)^{g^{\frac{n-1}{2}}},
\]
or also

\[
\sqrt[n]{R_{\nu+\frac{n-1}{2}}}\ \sqrt[n]{R_\nu} = \Phi(k_\nu) \left(\sqrt[n]{R_\nu}\right)^{g^{\frac{n-1}{2}}+1}.
\]

Now because $g$ is a primitive root of $n$, $g^{\frac{n-1}{2}}+1$ is divisible by $n$, and the right side therefore becomes rational. If then $\Psi$ denotes a rational function, we have

\[
\sqrt[n]{R_{\nu+\frac{n-1}{2}}}\ \sqrt[n]{R_\nu} = \Psi(k_\nu),
\]
and this formula shows that

\[
\Psi(k_\nu) = \Psi\!\left(k_{\nu+\frac{n-1}{2}}\right),
\]
hence that $\Psi(k_\nu)$ is real. From this it follows that

\[
\sqrt[n]{R_{\nu+\frac{n-1}{2}}} \quad \text{and} \quad \sqrt[n]{R_\nu}
\]
%CONT>>
